% !TEX root = main.tex

This work implemented and compared four classifiers for 3D mid-air gesture
recognition on two domains from the dataset of Huang et al.~\cite{Huang2019}:
Dynamic Time Warping (DTW), Edit-distance on quantized symbolic sequences,
the Three-Cents ($3\mathbb{c}$) recognizer, and a Bidirectional LSTM.
The three research questions raised in the introduction can now be answered.

\paragraph{Q1 --- Classical vs.\ modern methods.}
Contrary to the common expectation that neural or more elaborate methods
should dominate, the Three-Cents recognizer --- a lightweight
resample--scale--translate template matcher --- achieves the highest
user-independent accuracy on both domains ($96.3\%$ on domain~1,
$95.1\%$ on domain~4), outperforming the BiLSTM by a margin of $11$--$25$
percentage points. This result highlights that the \emph{inductive bias} of
the method matters more than its computational complexity: Three-Cents
implicitly encodes the relevant invariances (translation, scale, partial
rotation) that mid-air gesture trajectories satisfy, whereas the BiLSTM must
learn these from a dataset too small to generalize reliably across users.
The dynamic-programming baselines (DTW, Edit-distance) occupy an intermediate
range, with DTW consistently outperforming Edit-distance in the
user-independent regime.

\paragraph{Q2 --- User-independent vs.\ user-dependent accuracy.}
The user-dependent setting yields systematically higher accuracy for every
method and every domain, as expected. The gap is most pronounced for
template-based methods: DTW and Edit-distance improve by up to $38$ percentage
points in the user-dependent regime, where training and test gestures share
the same motor style. Three-Cents and, to a lesser extent, the BiLSTM are
more robust to user shift: their user-dependent gains are smaller precisely
because their user-independent performance is already high. In user-dependent
mode, DTW, Edit-distance, and Three-Cents all approach saturation
($\geq 98\%$), confirming that the task becomes near-trivial once
within-user templates are available.

\paragraph{Q3 --- Is one method statistically superior?}
Yes. The Friedman test rejects the null hypothesis of equal expected ranks in
both domains ($p < 0.001$), and Nemenyi and Wilcoxon--Holm post-hoc tests
consistently identify Three-Cents as the sole method significantly better than
all others in the user-independent setting. No other pairwise ordering is
consistent across both domains.

\paragraph{Limitations and perspectives.}
The BiLSTM's underperformance is most plausibly a data-size effect: with
$\sim$1\,000 samples per domain, the network cannot learn the cross-user
invariances that Three-Cents enforces by design. Training on a larger,
more diverse cohort --- or augmenting the dataset with synthetic inter-user
style transfers --- would likely narrow this gap considerably. The ablation
study further revealed that PCA applied without prior normalization can
degrade performance for Edit-distance, a methodological pitfall worth
flagging for practitioners. Finally, the confusion analysis showed that
all remaining errors are geometrically motivated: the Pyramid--Cylindrical
Pipe and digit~0--8 boundaries are the hardest pairs and could be targeted
with class-specific augmentation or cost-sensitive training.

More broadly, this project illustrates that well-designed, geometry-aware
feature engineering can rival or surpass end-to-end learned representations
when labeled data is scarce --- a conclusion that remains practically relevant
beyond the specific gesture recognition setting studied here.
